% !TeX program = xelatex
% ============================================================================
%  第 2 课　从余数到模运算：单位元与逆元　—— 教师版讲义
%  与 02-课件/第02课-从余数到模运算/第02课-从余数到模运算.tex 同步
% ============================================================================

\jhlesson{2}{从余数到模运算：单位元与逆元}{时长 45 分钟\quad·\quad 教具：12 格圆盘单页（人手一张）\quad·\quad 本课发单页}

\begin{mubiao}
  \begin{itemize}
    \item 会取余数，会写同余的记号 \(a\equiv b\pmod n\)。
    \item 能找到模 \(12\) 加法的单位元 \(0\)，以及 \(0,1,\dots,11\) 每个元素的逆元。
    \item 能说出有理数加法、有理数乘法里对应的单位元和逆元。
  \end{itemize}
\end{mubiao}

\noindent{\small 本课节奏：0—8 取余数｜8—14 模记号｜14—30 模 12 加法与单位元｜30—40 逆元｜40—45 和有理数对上。全程至少 4 次「用手指回答」。}

\section*{一、从余数说起（0—8 分钟）}

开场先算一笔：\textbf{今天是星期三，100 天以后是星期几？}

\begin{caozuo}
  转身板书 \(100=7\times14+2\)。不要急着说答案，先让全班看这个式子。
\end{caozuo}

\(7\) 天一循环，\(14\) 个整星期可以整块扔掉，只剩下 \(2\) 天，所以是\textbf{星期五}。

带余除法他们早就会了：

\[
  \underbrace{100}_{\text{被除数}} = \underbrace{7}_{\text{除数}}\times\underbrace{14}_{\text{商}} + \underbrace{2}_{\text{余数}} .
\]

以前关心的都是\textbf{商}；从今天起，只关心\textbf{余数}。

\section*{二、给余数一个记号（8—14 分钟）}

板书：
\[
  100\equiv2 \pmod 7 .
\]
读作「\(100\) 与 \(2\) 对模 \(7\) 同余」。

\begin{definition}[同余]
  如果两个整数 \(a,b\) 除以 \(n\) 之后的余数相同，就记作
  \[
    a\equiv b \pmod n,
  \]
  读作「\(a\) 与 \(b\) 对模 \(n\) 同余」。
\end{definition}

这句话翻成大白话就是：\(100\) 除以 \(7\) 的余数是 \(2\)。

\begin{caozuo}
  马上练两个，请全班用手指数表示余数：\(365\pmod 7\)（答案 \(1\)）、
  \(1000\pmod 7\)（答案 \(6\)）。这两个余数，就是「再过 365 天是星期几」
  「再过 1000 天是星期几」。
\end{caozuo}

\begin{zhu}
  这块点到为止。学生对带余除法非常熟，不要在这里磨时间。
\end{zhu}

\section*{三、模 12 的加法与单位元（14—30 分钟）}

黑板上画一个 \(12\) 格圆盘：\(0,1,2,\dots,11\) 顺时针排一圈，走满 \(12\) 格回到原地。
（单页上印的就是这个圆盘，学生人手一张。）

当场演示两个：

\begin{example}
  \(8+7=15\)。走满一圈就回到原点，所以 \(8+7\equiv3\pmod{12}\)。
\end{example}

\begin{example}
  \(11+1=12\equiv0\pmod{12}\)——从 \(11\) 再走一格，正好回到 \(0\)。
\end{example}

\begin{caozuo}
  手指回答，报出下列各式的答案（模 \(12\)）：\(10+5\)（\(3\)）、\(7+7\)（\(2\)）、
  \(9+4\)（\(1\)）、\(6+6\)（\(0\)）、\(3+3+3+3+3\)（\(3\)）。
\end{caozuo}

然后抛问题：在这套运算里，哪一个数 \(e\) 满足「加进去等于没加」，
也就是对每个 \(a\) 都有 \(a+e=a\)？——答案是 \(0\)，走 \(0\) 格就是不动。

\begin{definition}[单位元]
  设 \(*\) 是集合 \(S\) 上的一个运算。如果存在 \(e\in S\)，使得
  \[
    a*e=e*a=a
  \]
  对集合里\textbf{每一个}元素 \(a\) 都成立，就称 \(e\) 是这种运算的\textbf{单位元}。
\end{definition}

\begin{zhu}
  「单位」就是「一个」的意思：它不改变任何东西。
\end{zhu}

\section*{四、逆元（30—40 分钟）}

问：每个数要走几格，才能回到 \(0\)？把 \(12\) 个数两两配起来：

\begin{center}
{\small
\begin{tabular}{@{}l*{12}{c}@{}}
  \toprule
  \(a\) & 0 & 1 & 2 & 3 & 4 & 5 & 6 & 7 & 8 & 9 & 10 & 11 \\
  \midrule
  搭档 & 0 & 11 & 10 & 9 & 8 & 7 & 6 & 5 & 4 & 3 & 2 & 1 \\
  \bottomrule
\end{tabular}}
\end{center}

\(1\) 配 \(11\)（\(1+11=12\equiv0\)），\(2\) 配 \(10\)，\(3\) 配 \(9\)，……

\begin{zhu}
  特别指出：有两个人只配自己——\(0\)（\(0+0=0\)）和 \(6\)（\(6+6=12\equiv0\)）。
\end{zhu}

\begin{definition}[逆元]
  设 \(e\) 是运算 \(*\) 的单位元。如果对元素 \(a\) 能找到一个元素 \(b\)，使得
  \[
    a*b=b*a=e,
  \]
  就称 \(b\) 是 \(a\) 的\textbf{逆元}，也说 \(a\) 与 \(b\) \textbf{互逆}。
\end{definition}

\begin{dingli}[模 12 加法中逆元的存在性]
  在模 \(12\) 的加法里，每个元素都有逆元，就是上表里的搭档。
\end{dingli}

\section*{五、和有理数对上（40—45 分钟）}

板书一张对照表：

\begin{center}
\begin{tabular}{@{}lcc@{}}
  \toprule
   & \textbf{单位元} & \textbf{逆元} \\
  \midrule
  有理数加法 & \(0\) & 相反数 \(-a\) \\
  有理数乘法 & \(1\) & 倒数 \(1/a\ (a\neq0)\) \\
  模 \(12\) 加法 & \(0\) & 搭档（\(a\) 配 \(12-a\)） \\
  \bottomrule
\end{tabular}
\end{center}

一句话：这两样东西他们早就在用了，只是以前没这么叫过。

\begin{zhu}
  负数是为了「加法能走回头路」才被造出来的：\(5\) 加上 \(-5\) 得到单位元 \(0\)。
\end{zhu}

\begin{xiaojie}
  \begin{itemize}
    \item 取余数是有名字的算术：模 \(n\) 的加法，记号是 \(a\equiv b\pmod n\)。
    \item 模 \(12\) 加法里，单位元是 \(0\)，每个元素都有逆元（搭档）。
    \item 有理数加法、有理数乘法、模 \(12\) 加法，都有单位元和逆元。
  \end{itemize}
\end{xiaojie}

\noindent\textbf{下节课}：换一个运算试试——乘法还能不能找到逆元？